\documentclass[12pt,a4paper]{article}

\usepackage[a4paper,margin=1in]{geometry}
\usepackage[T1]{fontenc}
\usepackage[utf8]{inputenc}
\usepackage{lmodern}
\usepackage{setspace}
\usepackage{graphicx}
\usepackage{booktabs}
\usepackage{array}
\usepackage{tabularx}
\usepackage{longtable}
\usepackage{enumitem}
\usepackage{capt-of}
\usepackage{csquotes}
\usepackage[hidelinks]{hyperref}
\usepackage[round,authoryear]{natbib}

\setstretch{1.15}
\setlength{\parindent}{0pt}
\setlength{\parskip}{0.7em}
\setlist[itemize]{leftmargin=1.5em}
\setlist[enumerate]{leftmargin=1.8em}
\newcommand{\studentname}{Yiwei LI}
\newcommand{\studentid}{20513831}
\newcommand{\modulecode}{COMP3052}
\newcommand{\reporttitle}{Lab Report - Passwords}
\newcommand{\submissiondate}{22 March 2026}
\newcommand{\reportwordcount}{958}

\newcommand{\appendixfigurecaption}[2]{%
	\captionof{figure}{#1}
	\label{#2}
	\vspace{0.6em}
}

\newcommand{\appendixwideimage}[2][0.9\linewidth]{%
	\includegraphics[width=#1]{figures/#2}\par\vspace{0.7em}
}

\newcommand{\appendiximagepair}[2]{%
	\noindent
	\includegraphics[width=0.49\linewidth]{figures/#1}\hfill
	\includegraphics[width=0.49\linewidth]{figures/#2}\par\vspace{0.7em}
}

\begin{document}

\hypersetup{pageanchor=false}
\begin{titlepage}
	\centering
	\vspace*{2.4cm}
	{\Large \modulecode \par}
	\vspace{0.7cm}
	{\huge \reporttitle \par}
	\vspace{1.4cm}
	{\large Coursework 01 \par}
	\vfill
		\begin{tabular}{>{\bfseries}p{4.2cm}p{8cm}}
			Name: & \studentname \\
			Student ID: & \studentid \\
			Module Code: & \modulecode \\
			Date: & \submissiondate \\
			Word Count: & \reportwordcount \\
		\end{tabular}
	\vfill
	{\large University of Nottingham Ningbo China \par}
\end{titlepage}
\clearpage
\hypersetup{pageanchor=true}

\section{Introduction}
Passwords have been widely used for authentication in modern networks services and computational systems in recent decades, whereas they are known to be vulnerable to various attacks. This report proposes a password and authentication policy for a network with both ordinary users and privileged administrators. The design combines
observations from the lab experience with current guidance from NIST and the UK National Cyber Security
Centre (NCSC).

\section{Proposed Password and Authentication Policy}
The following policy for the network is proposed as:

\begin{itemize}
	\item All user passwords must be at least 15 characters long. 
	\item Administrator or sudo-capable accounts should not reuse passwords on any other account.
	\item Password validity should be limited to 30 days, with a warning period of 3 days and a
	further 3-day grace period before inactivation.
	\item Passwords must contain upper-case, lower-case, numeric, and special characters, in line
	with the stronger PAM checks explored in the lab.
	\item The system must block common passwords, dictionary words, leaked passwords, and obvious variants of those terms.
	\item High-risk accounts, especially administrators, remote access, VPN, and internet-facing
	accounts, must use multi-factor authentication (MFA).
	\item Passwords must never be stored in plain text. They should be protected using a salted,
	adaptive password hashing scheme.
	\item The login service should temporarily lock an account after five failed
	logins, and all logon attempts should be logged for audition.
\end{itemize}

The policy aligns with the lab materials in terms of PAM enforcement, password aging, and
format checks configuration. It aims to enhance password security and reduce the risk of being cracked when facing malicious attacks.
\section{Objective Justification}
This policy is supported by lab evidence and modern guidance. NIST states that single-factor passwords should be at least 15 characters
long, that verifiers should allow passwords of at least 64 characters, compare proposed
passwords against blocklists of common and compromised values, and avoid relying only on
simple composition formulae \citep{nist80063b2025}. The policy therefore adopts a 15-character
minimum, paired with character-mix checks, deny lists, password history, rate limiting, and
MFA for privileged access.

NIST also states that stored passwords should be salted and hashed using a tunable scheme
\citep{nist80063b2025}. NCSC recommends lockout and MFA for important accounts
\citep{ncscPolicy2018}, and notes that long multi-word passwords can improve usability
\citep{ncscThreeWords2021}. Lab evidence supports this. Appendix
Figure~\ref{fig:appendix-default-pam-checks} shows the default \texttt{pam\_unix}
configuration and a weak-password change attempt,
Figure~\ref{fig:appendix-strengthened-pam-checks} records the stricter
\texttt{pam\_cracklib} rule, and Figure~\ref{fig:appendix-hash-and-cracking-demo} in Appendix
shows that weak passwords remain recoverable after hashing. Figure~\ref{fig:appendix-user-creation-sudo} shows why sudo-capable accounts deserve stronger
protection, while Figure~\ref{fig:appendix-chage-policy} confirms that Linux can
enforce the intended 30/3/3 aging settings. The additional \texttt{hashcat} comparisons in Figure~\ref{fig:appendix-rule-attack-comparison} and
Figure~\ref{fig:appendix-mask-attack-comparison} show that some passwords yield to rule-based
transformations, while others depend on short searchable structure.

It is found that password-composition policies can impose usability costs without commensurate security gains \citep{egelman2011passwords}.
\citet{florencio2007strong} argue that, in online systems with limited login attempts, the
security return from increasingly elaborate composition rules is often modest once a reasonable
baseline has been met. For higher-value accounts, stronger protection is better achieved by
additional authentication. \citet{doerfler2019login} report that device-based login challenges
blocked more than 94\% of phishing-rooted hijacking attempts in their study, while
\citet{thomas2019breach} claimed that 1.5\% of observed web logins involved breached
credentials. Taken together, these findings justify deny lists, rate limiting, and MFA as
necessary complements to password rules.

\section{Reflection on Lab Experience}
The tasks in labs equipped us with hands-on experiences on hashcat tool for password cracking and analysis use, throwing insights on the design of strong and hard-to-crack passwords building upon lessons in historical password leakage.

The contrast between the default behaviour in
Figure~\ref{fig:appendix-default-pam-checks} and the stricter configuration in
Figure~\ref{fig:appendix-strengthened-pam-checks} showed that basic rejection (in traditional pam) is not the same
as strong enforcement. The default configuration offered only limited resistance to patterned
choices, so it was difficult to defend a policy based on traditional Linux settings alone.

In addition, Figure~\ref{fig:appendix-chage-policy} demonstrates that
password-age settings can be changed quickly with \texttt{chage}, while Appendix
Figure~\ref{fig:appendix-user-creation-sudo} shows that \texttt{sudo} privileges immediately
change the risk profile of an account. Password policy therefore cannot be separated from
privilege management in modern network and systems.

In conclusion, the lab experience demonstrates that a defensible password policy should reduce reliance on user memory tricks and shift security
work toward deny lists, rate limiting, secure hash storage, monitoring.

\section{Additional Experiments}
Going beyond the tasks in lab sessions, I also have used the \texttt{hashcat} tool to compare attack
strategies rather than simply repeating the supplied examples. I began The first experiment,
shown in Figure~\ref{fig:appendix-rule-attack-comparison}, by resetting the provided hash
set, running a straight dictionary attack with \texttt{dicts/small.txt}, and then repeated the test
with \texttt{rules/best64.rule} applied to the same dictionary. In my attempt, the straight attack
recovered only one password, whereas the rule-based version recovered 53. I inspected that the produced result is
hugely different because the extra recoveries came from 64 patterns of predictable edits that users often treat as
improvement, such as case changes, suffixes, or minor transformations. It therefore verifies 
that simple complexity tricks are not enough when an attacker applies widely known standard
cracking rules. 

I started another experiment as shown in Figure~\ref{fig:appendix-mask-attack-comparison}. It uses mask attacks to test whether short
structured passwords can be recovered even when they are not ordinary dictionary words. I
compared an upper-case seven-character mask with a short all-character mask and then
reviewed the recovered entries. In my run, both masks exhausted their search spaces quickly in minutes, as the length and rules were rather computationally effective,
but produced no new recoveries beyond the passwords already present in the potfile. This is
still informative because it shows that mask attacks are highly dependent on selecting the
right pattern in advance, indicating that this approach could be less-effective to handle certain passwords with limited compositional knowledge. 

Together, these experiments extend the basic lab knowledge by linking the
proposed policy to two realistic attack modes in hashcat: rule-based guessing and structured masked brute-force search.



%TC:ignore
\bibliographystyle{plainnat}
\bibliography{references}

\appendix

\section{Appendix: Evidence}

\subsection{Core Lab Evidence}
\appendixfigurecaption{Creation of the test account \texttt{uriforcw1} and group information modification.}{fig:appendix-user-creation-sudo}
\begingroup
\centering
\appendixwideimage[0.72\linewidth]{appendix-user-creation-sudo-panel-a-b.png}
\appendiximagepair{appendix-user-creation-sudo-panel-c.png}{appendix-user-creation-sudo-panel-d.png}
\endgroup

\appendixfigurecaption{Password aging policy modification.}{fig:appendix-chage-policy}
\begingroup
\centering
\appendixwideimage[0.82\linewidth]{appendix-chage-policy.png}
\endgroup

\appendixfigurecaption{Default PAM checks.}{fig:appendix-default-pam-checks}
\begingroup
\centering
\appendixwideimage[0.94\linewidth]{appendix-default-pam-checks.png}
\endgroup

\appendixfigurecaption{Strengthened PAM checks.}{fig:appendix-strengthened-pam-checks}
\begingroup
\centering
\appendixwideimage[0.94\linewidth]{appendix-strengthened-pam-checks-panel-a-b.png}
\appendixwideimage[0.94\linewidth]{appendix-strengthened-pam-checks-panel-c.png}
\appendixwideimage[0.94\linewidth]{appendix-strengthened-pam-checks-panel-d-e.png}
\endgroup

\appendixfigurecaption{Hash and cracking demo.}{fig:appendix-hash-and-cracking-demo}
\begingroup
\centering
\appendixwideimage[0.9\linewidth]{appendix-hash-and-cracking-demo-panel-a-b-c-d.png}
\appendixwideimage[0.82\linewidth]{appendix-hash-and-cracking-demo-panel-e-f.png}
\endgroup

\subsection{Additional Experiment 1: Straight Attack Versus Rule Attack}
\appendixfigurecaption{Straight attack and rule attack comparison.}{fig:appendix-rule-attack-comparison}
\begingroup
\centering
\appendixwideimage[0.86\linewidth]{appendix-rule-attack-comparison-panel-a-b.png}
\appendixwideimage[0.78\linewidth]{appendix-rule-attack-comparison-panel-d.png}
\appendixwideimage[0.86\linewidth]{appendix-rule-attack-comparison-panel-e.png}
\appendixwideimage[0.78\linewidth]{appendix-rule-attack-comparison-panel-f.png}
\endgroup

\subsection{Additional Experiment 2: Targeted Mask Attacks}
\appendixfigurecaption{Targeted mask attacks.}{fig:appendix-mask-attack-comparison}
\begingroup
\centering
\appendixwideimage[0.86\linewidth]{appendix-mask-attack-comparison-panel-a-b-c.png}
\appendixwideimage[0.78\linewidth]{appendix-mask-attack-comparison-panel-d.png}
\appendixwideimage[0.86\linewidth]{appendix-mask-attack-comparison-panel-e.png}
\appendixwideimage[0.78\linewidth]{appendix-mask-attack-comparison-panel-f-h.png}
\endgroup
%TC:endignore

\end{document}
